%%%%%%%%%%%%%%%%
%%% deut 21 %%%
%%%%%%%%%%%%%%%%

\Chapter{21}

\Verse{1}
{
	\hDtn{כִּי יִמָּצֵא חָלָל בָּאֲדָמָה אֲשֶׁר יְהֹוָה אֱלֹהֶיךָ נֹתֵן לְךָ לְרִשְׁתָּהּ נֹפֵל בַּשָּׂדֶה לֹא נוֹדַע מִי הִכָּהוּ׃}
}
{
	\eDtn{
		21 “If a corpse will be found in the land that YHWH, your
		{\God}, is giving you to take possession of it, fallen in the
		field, unknown who struck him,
	}
}
{
}

\Verse{2}
{
	\hDtn{וְיָצְאוּ זְקֵנֶיךָ וְשֹׁפְטֶיךָ וּמָדְדוּ אֶל הֶעָרִים אֲשֶׁר סְבִיבֹת הֶחָלָל׃}
}
{
	\eDtn{
		then your elders and your judges shall go out and measure
		to the cities that are around the corpse.
	}
}
{
}

\Verse{3}
{
	\hDtn{וְהָיָה הָעִיר הַקְּרֹבָה אֶל הֶחָלָל וְלָקְחוּ זִקְנֵי הָעִיר הַהִוא עֶגְלַת בָּקָר אֲשֶׁר לֹא עֻבַּד בָּהּ אֲשֶׁר לֹא מָשְׁכָה בְּעֹל׃}
}
{
	\eDtn{
		And it will be, the closest city to the corpse: that
		city’s elders shall take a heifer that has not been worked
		with, that has not pulled in a yoke.
	}
}
{
}

\Verse{4}
{
	\hDtn{וְהוֹרִדוּ זִקְנֵי הָעִיר הַהִוא אֶת הָעֶגְלָה אֶל נַחַל אֵיתָן אֲשֶׁר לֹא יֵעָבֵד בּוֹ וְלֹא יִזָּרֵעַ וְעָרְפוּ שָׁם אֶת הָעֶגְלָה בַּנָּחַל׃}
}
{
	\eDtn{
		And that city’s elders shall take the heifer down to a
		strongly flowing wadi that would not be worked with and
		would not be seeded, and they shall break the heifer’s
		neck there in the wadi.
	}
}
{
}

\Verse{5}
{
	\hDtn{וְנִגְּשׁוּ הַכֹּהֲנִים בְּנֵי לֵוִי כִּי בָם בָּחַר יְהֹוָה אֱלֹהֶיךָ לְשָׁרְתוֹ וּלְבָרֵךְ בְּשֵׁם יְהֹוָה וְעַל פִּיהֶם יִהְיֶה כּל רִיב וְכל נָגַע׃}
}
{
	\eDtn{
		And the priests, sons of Levi, shall go over, because
		YHWH, your {\God}, chose them to serve Him and to bless in
		YHWH’s name, and every dispute and every injury shall be
		by their word.
	}
}
{
}

\Verse{6}
{
	\hDtn{וְכֹל זִקְנֵי הָעִיר הַהִוא הַקְּרֹבִים אֶל הֶחָלָל יִרְחֲצוּ אֶת יְדֵיהֶם עַל הָעֶגְלָה הָעֲרוּפָה בַנָּחַל׃}
}
{
	\eDtn{
		And all of that city’s elders, who are close to the
		corpse, shall wash their hands over the heifer whose neck
		was broken in the wadi.
	}
}
{
}

\Verse{7}
{
	\hDtn{וְעָנוּ וְאָמְרוּ יָדֵינוּ לֹא (שפכה) [שָׁפְכוּ] אֶת הַדָּם הַזֶּה וְעֵינֵינוּ לֹא רָאוּ׃}
}
{
	\eDtn{
		And they shall testify, and they shall say, ‘Our hands did
		not spill this blood, and our eyes did not see it.
	}
}
{
}

\Verse{8}
{
	\hDtn{כַּפֵּר לְעַמְּךָ יִשְׂרָאֵל אֲשֶׁר פָּדִיתָ יְהֹוָה וְאַל תִּתֵּן דָּם נָקִי בְּקֶרֶב עַמְּךָ יִשְׂרָאֵל וְנִכַּפֵּר לָהֶם הַדָּם׃}
}
{
	\eDtn{
		Grant atonement for your people Israel, whom you redeemed,
		YHWH, and don’t impute innocent blood among your people
		Israel.’ And for them the blood will be atoned for.
	}
}
{
}

\Verse{9}
{
	\hDtn{וְאַתָּה תְּבַעֵר הַדָּם הַנָּקִי מִקִּרְבֶּךָ כִּי תַעֲשֶׂה הַיָּשָׁר בְּעֵינֵי יְהֹוָה׃}
}
{
	\eDtn{
		So you shall burn away the innocent blood from among you
		when you will do what is right in YHWH’s eyes.
	}
}
{
}

\Verse{10}
{
	\hDtn{כִּי תֵצֵא לַמִּלְחָמָה עַל אֹיְבֶיךָ וּנְתָנוֹ יְהֹוָה אֱלֹהֶיךָ בְּיָדֶךָ וְשָׁבִיתָ שִׁבְיוֹ׃}
}
{
	\eDtn{
		WHEN YOU’LL GO OUT “When you’ll go out to war against your
		enemies, and YHWH, your {\God}, will put him in your hand,
		and you’ll take prisoners from him,
	}
}
{
}

\Verse{11}
{
	\hDtn{וְרָאִיתָ בַּשִּׁבְיָה אֵשֶׁת יְפַת תֹּאַר וְחָשַׁקְתָּ בָהּ וְלָקַחְתָּ לְךָ לְאִשָּׁה׃}
}
{
	\eDtn{
		and you’ll see among the prisoners a woman with a
		beautiful figure, and you’ll be attracted to her and take
		her for yourself as a wife,
	}
}
{
%	\footnote{
%		a woman among the prisoners. A number of this law’s details have never
%		been explained definitively. Some would say she cuts her hair and nails
%		so as to be less attractive, so that the man who is attracted to her
%		will become less infatuated. Some would say it is part of her mourning.
%		But what pervades the elements of the law is an extraordinary
%		sensitivity to the humanity of a captured woman. Contrary to one of the
%		most common practices of war, the Israelite soldier is not permitted to
%		rape her. He may take her as a wife. But even then he must give her time
%		to mourn the loss of her family. And if he takes her as a wife but then
%		rejects her, he must let her go completely free. The text uses the same
%		verb as in the law of divorce (Deut 24:1). And the text recognizes that
%		he has degraded her, and so he cannot treat her like a slave to be sold;
%		he cannot receive money for her from anyone else. Notably, the words for
%		degrading her and for letting her go are the same words that are used to
%		describe the Egyptians’ degrading of Israel and then letting Israel go
%		(Exod 1:11–12; 5:1). Again Israel learns from its experience of
%		enslavement. Israel not only celebrates its own release, but it learns
%		to have compassion for others as well.
%	}
}

\Verse{12}
{
	\hDtn{וַהֲבֵאתָהּ אֶל תּוֹךְ בֵּיתֶךָ וְגִלְּחָה אֶת רֹאשָׁהּ וְעָשְׂתָה אֶת צִפּרְנֶיהָ׃}
}
{
	\eDtn{
		then you shall bring her into your house, and she shall
		shave her head and do her nails
	}
}
{
}

\Verse{13}
{
	\hDtn{וְהֵסִירָה אֶת שִׂמְלַת שִׁבְיָהּ מֵעָלֶיהָ וְיָשְׁבָה בְּבֵיתֶךָ וּבָכְתָה אֶת אָבִיהָ וְאֶת אִמָּהּ יֶרַח יָמִים וְאַחַר כֵּן תָּבוֹא אֵלֶיהָ וּבְעַלְתָּהּ וְהָיְתָה לְךָ לְאִשָּׁה׃}
}
{
	\eDtn{
		and take away her prisoner’s garment from on her. And she
		shall live in your house and shall mourn her father and
		her mother a month of days. And after that you may come to
		her and marry her, and she shall become your wife.
	}
}
{
}

\Verse{14}
{
	\hDtn{וְהָיָה אִם לֹא חָפַצְתָּ בָּהּ וְשִׁלַּחְתָּהּ לְנַפְשָׁהּ וּמָכֹר לֹא תִמְכְּרֶנָּה בַּכָּסֶף לֹא תִתְעַמֵּר בָּהּ תַּחַת אֲשֶׁר עִנִּיתָהּ׃}
}
{
	\eDtn{
		And it will be, if you don’t desire her then you shall let
		her go on her own, and you shall not sell her for money.
		You shall not get profit through her, because you degraded
		her.
	}
}
{
}

\Verse{15}
{
	\hDtn{כִּי תִהְיֶיןָ לְאִישׁ שְׁתֵּי נָשִׁים הָאַחַת אֲהוּבָה וְהָאַחַת שְׂנוּאָה וְיָלְדוּ לוֹ בָנִים הָאֲהוּבָה וְהַשְּׂנוּאָה וְהָיָה הַבֵּן הַבְּכֹר לַשְּׂנִיאָה׃}
}
{
	\eDtn{
		“When a man will have two wives, one loved and one hated,
		and they’ll give birth to children for him, the loved and
		the hated, and the hated will have the firstborn son,
	}
}
{
}

\Verse{16}
{
	\hDtn{וְהָיָה בְּיוֹם הַנְחִילוֹ אֶת בָּנָיו אֵת אֲשֶׁר יִהְיֶה לוֹ לֹא יוּכַל לְבַכֵּר אֶת בֶּן הָאֲהוּבָה עַל פְּנֵי בֶן הַשְּׂנוּאָה הַבְּכֹר׃}
}
{
	\eDtn{
		then it will be, on the day that he gives what he has as a
		legacy to his children, he shall not be able to give a son
		of the loved one the birthright before the firstborn son
		of the hated one.
	}
}
{
}

\Verse{17}
{
	\hDtn{כִּי אֶת הַבְּכֹר בֶּן הַשְּׂנוּאָה יַכִּיר לָתֶת לוֹ פִּי שְׁנַיִם בְּכֹל אֲשֶׁר יִמָּצֵא לוֹ כִּי הוּא רֵאשִׁית אֹנוֹ לוֹ מִשְׁפַּט הַבְּכֹרָה׃}
}
{
	\eDtn{
		But he shall recognize the firstborn son of the hated one,
		to give him a double portion of all that he has, because
		he is the beginning of his might. The legal due of the
		firstborn is his.
	}
}
{
%	\footnote{
%		beginning of his might. Moses uses the same phrase that Jacob uses to
%		describe Reuben, his firstborn, when he speaks of his sons’ legacies
%		(Gen 49:3). Perhaps this was a known expression for firstborn sons, or
%		perhaps we should picture Moses as deliberately bringing Jacob’s words
%		to mind. It is ironic because Jacob does in fact take away the double
%		portion from Reuben, who is the son of the less-loved wife, Leah; and he
%		gives it to Joseph, who is the son of the more-loved wife, Rachel (Gen
%		29:30). It may be that Moses is teaching that the Israelites may not do
%		what Jacob did. Or it may be that Jacob’s case demonstrates the point of
%		the law that Moses is giving, since Jacob does not take the birthright
%		away from Reuben because of Reuben’s mother, Leah, being less loved. He
%		takes it away because Reuben himself committed an offense, sleeping with
%		his father’s concubine (Gen 35:22; 49:4).
%	}
}

\Verse{18}
{
	\hDtn{כִּי יִהְיֶה לְאִישׁ בֵּן סוֹרֵר וּמוֹרֶה אֵינֶנּוּ שֹׁמֵעַ בְּקוֹל אָבִיו וּבְקוֹל אִמּוֹ וְיִסְּרוּ אֹתוֹ וְלֹא יִשְׁמַע אֲלֵיהֶם׃}
}
{
	\eDtn{
		“When a man will have a stubborn and rebellious son, not
		listening to his father’s voice and his mother’s voice,
		and they will discipline him, but he will not listen to
		them,
	}
}
{
}

\Verse{19}
{
	\hDtn{וְתָפְשׂוּ בוֹ אָבִיו וְאִמּוֹ וְהוֹצִיאוּ אֹתוֹ אֶל זִקְנֵי עִירוֹ וְאֶל שַׁעַר מְקֹמוֹ׃}
}
{
	\eDtn{
		then his father and his mother shall take hold of him and
		bring him out to his city’s elders and to the gate of his
		place.
	}
}
{
}

\Verse{20}
{
	\hDtn{וְאָמְרוּ אֶל זִקְנֵי עִירוֹ בְּנֵנוּ זֶה סוֹרֵר וּמֹרֶה אֵינֶנּוּ שֹׁמֵעַ בְּקֹלֵנוּ זוֹלֵל וְסֹבֵא׃}
}
{
	\eDtn{
		And they shall say to his city’s elders, ‘This son of ours
		is stubborn and rebellious, he doesn’t listen to our
		voice, a glutton, and a drunk.’
	}
}
{
}

\Verse{21}
{
	\hDtn{וּרְגָמֻהוּ כּל אַנְשֵׁי עִירוֹ בָאֲבָנִים וָמֵת וּבִעַרְתָּ הָרָע מִקִּרְבֶּךָ וְכל יִשְׂרָאֵל יִשְׁמְעוּ וְיִרָאוּ׃}
}
{
	\eDtn{
		And all his city’s people shall batter him with stones so
		he dies. So you shall burn away what is bad from among
		you. And all Israel will listen and fear.
	}
}
{
}

\Verse{22}
{
	\hDtn{וְכִי יִהְיֶה בְאִישׁ חֵטְא מִשְׁפַּט מָוֶת וְהוּמָת וְתָלִיתָ אֹתוֹ עַל עֵץ׃}
}
{
	\eDtn{
		“And if there will be a sin bringing a sentence of death
		on a man, and he will be put to death, and you will hang
		him on a tree,
	}
}
{
}

\Verse{23}
{
	\hDtn{לֹא תָלִין נִבְלָתוֹ עַל הָעֵץ כִּי קָבוֹר תִּקְבְּרֶנּוּ בַּיּוֹם הַהוּא כִּי קִלְלַת אֱלֹהִים תָּלוּי וְלֹא תְטַמֵּא אֶת אַדְמָתְךָ אֲשֶׁר יְהֹוָה אֱלֹהֶיךָ נֹתֵן לְךָ נַחֲלָה׃}
}
{
	\eDtn{
		you shall not leave his corpse on the tree, but you shall
		bury him on that day, because a hanged person is an
		offense to {\God}, and you shall not make impure your land
		that YHWH, your {\God}, is giving you as a legacy.
	}
}
{
}
